Un espectròmetre de masses és un aparell que permet determinar la relació càrrega/massa d'ions. L'espectròmetre de masses conté tres parts diferenciades. La primera part és un filament que ionitza les molècules o àtoms que entren dins l'espectròmetre. A la sortida del filament tots els ions tenen una càrrega negativa. A la segona part de l'aparell els ions passen per un selector de velocitats (figura 1) que està format per dues plaques paral·leles, entre les quals es genera un camp elèctric uniforme. La separació entre aquestes plaques és d'1,50~cm. Entre les plaques també es genera un camp magnètic uniforme de 0,50~T perpendicular al pla del paper i en sentit sortint, tal com es mostra en la figura 1.

\figura[0.4]{fig1.png}
\begin{center}
\textsc{Figura 1}
\end{center}

\begin{apartat}{1,25}
Volem que el selector de velocitats només deixi passar els ions que es moguin a una velocitat de $2,00\times 10^{5}\ \mathrm{m\,s^{-1}}$. Determineu la diferència de potencial que hem d'aplicar entre les plaques perquè els ions que es mouen a aquesta velocitat no es desviïn. Quina placa s'ha de connectar a potencial alt i quina a potencial baix? Justifiqueu les respostes i representeu les forces que actuen sobre un ió. Digueu si el selector de velocitats configurat d'aquesta manera també funciona per a ions positius i justifiqueu la resposta.
\end{apartat}

\figura[0.2]{fig2.png}
\begin{center}
\textsc{Figura 2}
\end{center}

\begin{apartat}{1,25}
La tercera part de l'espectròmetre es troba a la sortida del selector de velocitats i és una regió on hi ha un altre camp magnètic uniforme de 0,20~T, perpendicular al pla del paper i en sentit entrant (figura 2). Les pantalles laterals permeten mesurar la posició a què impacten els ions i d'aquesta manera poder determinar-ne la massa.\\
\hspace*{1.5em}Representeu esquemàticament sobre la figura 2 la trajectòria que descriuen els ions que surten del selector de velocitats indicant la direcció i el sentit de la força que exerceix el camp magnètic en un punt de la trajectòria. Justifiqueu la resposta. Calculeu a quina distància de la sortida del selector de velocitats impactarà l'ió dels isòtops del neó $^{20}\mathrm{Ne}^{-}$ (l'ió té la mateixa càrrega que un electró).
\end{apartat}

\dades{%
  $|e| = 1,602\times 10^{-19}\ \mathrm{C}$.\\
  Massa de l'ió: $^{20}\mathrm{Ne}^{-} = 3,32\times 10^{-26}\ \mathrm{kg}$.}
